% Part:sets-functions-relations
% Chapter: size-of-sets
% Section: zig-zag

\documentclass[../../../include/open-logic-section]{subfiles}

\begin{document}

\olfileid[fa-IR]{sfr}{siz}{zigzag}

\olsection{روش زیگزاگ کانتور}

\begin{explain}
پیش‌تر چند برشماریِ ``آسان'' را بررسی کرده‌ایم. اکنون موردی کمی
دشوارتر را بررسی می‌کنیم. مجموعهٔ زوج‌های مرتب اعداد طبیعی را در نظر
بگیرید\oliflabeldef{sfr:set:pai:sec}{، که آن را در
\olref[set][pai]{sec} چنین تعریف کردیم:}{، که چنین تعریف می‌شود:}
\[
\Nat \times \Nat = \Setabs{\tuple{n,m}}{n,m \in \Nat}
\]
می‌توانیم این زوج‌های مرتب را در یک \emph{آرایه} به‌شکل زیر سامان دهیم:
\[
\begin{array}{ c | c | c | c | c | c}
& \mathbf 0 & \mathbf 1 & \mathbf 2 & \mathbf 3 & \dots \\
\hline
\mathbf 0 & \tuple{0,0} & \tuple{0,1} & \tuple{0,2} & \tuple{0,3} & \dots \\
\hline
\mathbf 1 & \tuple{1,0} & \tuple{1,1} & \tuple{1,2} & \tuple{1,3} & \dots \\
\hline
\mathbf 2 & \tuple{2,0} & \tuple{2,1} & \tuple{2,2} & \tuple{2,3} & \dots \\
\hline
\mathbf 3 & \tuple{3,0} & \tuple{3,1} & \tuple{3,2} & \tuple{3,3} & \dots \\
\hline
\vdots & \vdots & \vdots & \vdots & \vdots & \ddots\\
\end{array}
\]
آشکار است که هر زوج مرتب در $\Nat \times \Nat$ دقیقاً یک‌بار در آرایه
ظاهر می‌شود. به‌ویژه، $\tuple{n,m}$ در سطر $n$اُم و ستون $m$اُم ظاهر
خواهد شد. اما چگونه اعضای چنین آرایه‌ای را در فهرستی ``یک‌بعدی'' سامان
دهیم؟ الگوی آرایهٔ زیر یک شیوهٔ انجام این کار را نشان می‌دهد (هرچند
البته گزینه‌های بسیار دیگری نیز وجود دارند):
\[
\begin{array}{ c | c | c | c | c | c | c}
& \mathbf 0 & \mathbf 1 & \mathbf 2 & \mathbf 3 & \mathbf 4 &\dots \\
\hline
\mathbf 0 & 0  & 1& 3 & 6& 10 &\ldots \\
\hline
\mathbf 1 &2 & 4& 7 & 11 & \dots &\ldots \\
\hline
\mathbf 2 & 5 & 8 & 12 & \ldots & \dots&\ldots \\
\hline
\mathbf 3 & 9 & 13 & \ldots & \ldots & \dots & \ldots \\
\hline
\mathbf 4 & 14 & \ldots & \ldots & \ldots & \dots & \ldots \\
\hline
\vdots & \vdots & \vdots & \vdots & \vdots&\ldots & \ddots\\
\end{array}
\]\noindent
این الگو \emph{روش زیگزاگ کانتور} نامیده می‌شود. این روش
$\Nat \times \Nat$ را به‌صورت زیر برمی‌شمارد:
\[
\tuple{0,0}, \tuple{0,1}, \tuple{1,0}, \tuple{0,2}, \tuple{1,1},
\tuple{2,0}, \tuple{0,3}, \tuple{1,2}, \tuple{2,1}, \tuple{3,0}, \dots
\]
و این، گزارهٔ زیر را ثابت می‌کند:
\end{explain}

\begin{prop}\ollabel{natsquaredenumerable}
$\Nat \times \Nat$ !!{enumerable} است.
\end{prop}

\begin{proof}
تابع $f \colon \Nat \to \Nat\times\Nat$ را چنان تعریف می‌کنیم که هر
$k \in \Nat$ را به زوج $\tuple{n,m} \in \Nat \times \Nat$ بفرستد،
به‌گونه‌ای که $k$
مقدار خانهٔ واقع در سطر $n$اُم و ستون $m$اُمِ آرایهٔ زیگزاگ کانتور باشد.
\end{proof}

\begin{explain}
این فن به‌خوبی قابل تعمیم نیز هست. برای نمونه، می‌توانیم از آن برای
برشمردن مجموعهٔ سه‌تایی‌های مرتب اعداد طبیعی استفاده کنیم؛ یعنی:
\[
\Nat \times \Nat \times \Nat = \Setabs{\tuple{n,m,k}}{n,m,k \in \Nat}
\]
$\Nat \times \Nat \times \Nat$ را ضرب دکارتیِ
$\Nat \times \Nat$ در $\Nat$ در نظر می‌گیریم؛ یعنی
\[
\Nat^3 = (\Nat \times \Nat) \times \Nat =
\Setabs{\tuple{\tuple{n,m},k}}{n, m, k
  \in \Nat }
\]
و ازاین‌رو می‌توانیم $\Nat^3$ را با آرایه‌ای برشماریم که یک محور آن
با برشماریِ $\Nat$ و محور دیگرش با برشماریِ $\Nat^2$ برچسب‌گذاری شده است:
\[
\begin{array}{ c | c | c | c | c | c}
& \mathbf 0 & \mathbf 1 & \mathbf 2 & \mathbf 3 & \dots \\
\hline
\mathbf{\tuple{0,0}} & \tuple{0,0,0} & \tuple{0,0,1} & \tuple{0,0,2} & \tuple{0,0,3} & \dots \\
\hline
\mathbf{\tuple{0,1}} & \tuple{0,1,0} & \tuple{0,1,1} & \tuple{0,1,2} & \tuple{0,1,3} & \dots \\
\hline
\mathbf{\tuple{1,0}} & \tuple{1,0,0} & \tuple{1,0,1} & \tuple{1,0,2} & \tuple{1,0,3} & \dots \\
\hline
\mathbf{\tuple{0,2}} & \tuple{0,2,0} & \tuple{0,2,1} & \tuple{0,2,2} & \tuple{0,2,3} & \dots\\
\hline
\vdots & \vdots & \vdots & \vdots & \vdots & \ddots \\
\end{array}
\]
پس با استفاده از روشی مانند روش زیگزاگ کانتور، می‌توانیم به همین ترتیب
برشماری‌ای از~$\Nat^3$ به دست آوریم. همچنین می‌توانیم این روند را ادامه
دهیم و برشماری‌هایی از $\Nat^n$ برای هر عدد طبیعی $n$ به دست آوریم.
بنابراین داریم:
\end{explain}

\begin{prop}
$\Nat^n$ برای هر $n \in \Nat$ !!{enumerable} است.
\end{prop}

\begin{prob}\label{sfr:siz:zigzag:prob:posint-n}
نشان دهید $(\PosInt)^n$ برای هر $n \in \Nat$ !!{enumerable} است.
\end{prob}

\begin{prob}\label{sfr:siz:zigzag:prob:posint-star}
  نشان دهید $(\PosInt)^*$ !!{enumerable} است. می‌توانید
  \cref{sfr:siz:zigzag:prob:posint-n} را فرض کنید.
\end{prob}

\end{document}
